\documentclass[11pt, a4paper]{article}
\usepackage[utf8]{inputenc}
\usepackage[T1]{fontenc}
\usepackage{amsmath, amssymb}
\usepackage{geometry}
\geometry{margin=1in}
\usepackage{booktabs}
\usepackage{hyperref}

\title{\textbf{ITALIAN OLYMPIADS IN INFORMATICS}\\Pre-EGOI\\Online, June 26, 2022}
\author{}
\date{}

\begin{document}
\maketitle
\vspace{-4em}

\section*{On the roofs of Pisa (parkour)}
\textbf{parkour IT}

Tommaso has recently become passionate about parkour. Unfortunately, Pisa, the city where Tommaso lives, does not have any gym where he can train, and therefore he is forced to practice on the roofs of the city.

The neighborhood where Tommaso trains consists of a row of $N$ houses, and the $i$-th ($0\le i<N$) is $S_{i}$ meters high.
When he is on roof $i$, Tommaso can make a jump of a length between $A_{i}$ and $B_{i}$ meters (inclusive).
Therefore, he can reach positions from $i+A_{i}$ to $i+B_{i}$. His goal is to reach the park, at position $N$.
Although he is a reckless boy, Tommaso is a victim of a condition quite adverse to the life of a traceur (parkour practitioner): he suffers from vertigo!
In fact, when he trains, he always tries to stay on the lowest roofs.
In particular, starting from roof $0$, Tommaso wants to find the path that takes him to position $N$ which minimizes the maximum height of the roofs he jumps on.
Help him in this task!

\subsection*{Implementation}
You will need to submit a single file, with a .cpp extension.
Among the attachments to this task, you will find a template \texttt{parkour.cpp} with an example implementation.
You will need to implement the following function:

\paragraph{C++}
\texttt{int salta(int N, std::vector<int> S, std::vector<int> A, std::vector<int> B);}
\begin{itemize}
    \item The integer $N$ represents the number of houses.
    \item The vector $S$, indexed from $0$ to $N-1$, contains the heights in meters of the houses.
    \item The vectors $A$ and $B$, indexed from $0$ to $N-1$, contain the bounds of the interval of houses onto which Tommaso can jump when he is on house $i$. In other words, Tommaso from the roof of house $i$ can jump to position $j$ with $A[i]\le j\le B[i]$.
\end{itemize}

The function must return the minimum possible maximum height of the roofs on a path that starts at house $0$ and arrives at position $N$.
The grader will call the \texttt{salta} function and will print its returned value to the output file.

\subsection*{Sample Grader}
In the directory relating to this problem there is a simplified version of the grader used during the evaluation, which you can use to test your solutions locally.
The sample grader reads data from \texttt{stdin}, calls the functions you need to implement, and writes to \texttt{stdout}, according to the following format.

The input file consists of $N+1$ lines, containing:
\begin{itemize}
    \item Line 1: the single integer $N$.
    \item Line $2+i$: the three values $S[i]$, $A[i]$, $B[i]$.
\end{itemize}

The output file consists of a single line, containing:
\begin{itemize}
    \item Line 1: the value returned by the \texttt{salta} function.
\end{itemize}

\subsection*{Assumptions}
\begin{itemize}
    \item $1\le N\le 1\,000\,000$.
    \item $0\le S[i]\le 1\,000\,000$ for each $i=0,\dots,N-1$.
    \item $1\le A[i]\le B[i]\le N-i$ for each $i=0,\dots,N-1$.
\end{itemize}

\subsection*{Scoring}
Your program will be tested on several test cases grouped into subtasks.
To obtain the score for a subtask, it is necessary to solve correctly all the test cases that comprise it.
\begin{itemize}
    \item Subtask 1 [0 points]: Sample cases.
    \item Subtask 2 [5 points]: $N\le 10$.
    \item Subtask 3 [10 points]: $N\le 100$ and $S[i]\le 100$ for each $i=0,\dots,N-1$.
    \item Subtask 4 [10 points]: $N\le 10\,000$ and $S[i]\le 10\,000$ for each $i=0,\dots,N-1$.
    \item Subtask 5 [15 points]: $N\le 10\,000$.
    \item Subtask 6 [30 points]: $N\le 100\,000$.
    \item Subtask 7 [10 points]: $A[i]=1$.
    \item Subtask 8 [20 points]: No additional limitations.
\end{itemize}

\subsection*{Input/Output Examples}

\begin{tabular}{@{}ll@{}}
\toprule
\textbf{stdin} & \textbf{stdout} \\
\midrule
4 & 2 \\
1 1 3 & \\
3 3 3 & \\
1 1 1 & \\
2 1 1 & \\
\midrule
6 & 2 \\
2 2 3 & \\
1 3 4 & \\
2 1 2 & \\
3 2 3 & \\
1 1 2 & \\
1 1 1 & \\
\bottomrule
\end{tabular}

\subsection*{Explanation}
In the first sample case, Tommaso starts from roof 0 and can reach roofs 1, 2, and 3. If he jumps onto roof 2, then onto roof 3, and finally arrives at the destination, the maximum height reached is 2. This is not the only solution, but it is not possible to reach the park by jumping only on roofs lower than 2.

In the second sample case, with the first jump he can reach roofs 2 and 3. If the path were 0, 2, 4, 6, the maximum height reached is 2. It is not possible to reach the park by jumping only on roofs lower than 2.

\end{document}
